\section{Audit, corrections, and verification}
\label{sec:audit}

The source audit was performed against the exact commit specified on
the title page. It was selective and mathematical: we checked the
definitions and proof dependencies needed for the four continuations,
then compared the conclusions with the manuscript's explicit open
questions. This is not an assertion that every formula in the entire
repository has been independently re-proved.

\subsection{A wording correction concerning non-elementarity}

The paragraph preceding Equation~\texttt{vertical:eq:Im3} in
\src{chapters/03-algebraic.tex} describes the imaginary part of
$\Li_3(1+i)$ as \emph{no longer elementary at all}. The displayed
hypergeometric representation does not establish that assertion.
An exact expression in terms of a central-binomial tower constant
is an evaluation in a specified language. It is not a proof that
the value cannot belong to some smaller algebra of constants.
Indeed, without specifying the meaning of elementary for a
number, the assertion has no definite mathematical test.

\begin{samepage}
A precise replacement is:
\begin{quote}
The imaginary part at $1+i$ is represented here by the
central-binomial Gaussian tower constant:
\end{quote}
\end{samepage}
followed by the existing, unchanged equation, and then:
\begin{quote}
This representation does not establish non-reducibility to a
specified smaller algebra of constants.
\end{quote}
The preceding claim that the general real part \emph{does not close}
should similarly be scoped to the computation being described.
One can say that differentiating and integrating the displayed
formula introduces an inverse-tangent-square integral, and that
no general reduction of that integral is supplied there.
This is a correction to the strength of the prose; the displayed
special-value formulas are not contradicted by this audit.
An exact-context patch implementing these two edits is included
in \src{integration/non_elementarity_wording.patch}.

\subsection{Statements that should be updated after integration}

\paragraph{Uniform Euler error.}
The warning that a bound on $|g_{a,b}|$ is insufficient to control
$2^N(E_N-g_{a,b})$ was correct. The explicit calculation in
Section~\ref{sec:euler} shows $R_2(1,1)>R_1(1,1)$.
Theorem~\ref{thm:Euler-uniform} now supplies the missing uniform
comparison. The open-status sentences in the subcritical and
real-order Euler sections should therefore be replaced by a
reference to the new theorem. The sharper constant $1$ on the
restricted domain $a\geq1$ remains useful and should be retained.
A global budget of $57/50$ is convenient, but it is not the sharp
bound on every subdomain.

\paragraph{Quartic threshold.}
The final paragraph of \src{chapters/05-real-turning.tex} correctly
labels its approximate simultaneous zero as diagnostic. Our exact
rectangle and derivative certificate promote that particular
location, its local uniqueness, and its sixth-order behavior to
theorems. They do not establish uniqueness on all $0<b<1$.
The current existence proof based on opposite signs of the quartic
coefficient remains a useful conceptual introduction. It should
be followed by the new local classification and explicit rational
two-extremum example, with a separate global open question.

\paragraph{Reflected coefficients and residues.}
The coefficient $A_{k,m}$ defined by Lagrange inversion and the
meromorphic residue are related by
\[
 \operatorname*{Res}_{z=k}
 \int_0^1\Gamma(x)^z\log^m\Gamma(1-x)\,dx=-kA_{k,m}.
\]
Any compact description of signs should retain this factor.
The moving-parameter theorems extend the fixed-$m$ estimates;
they do not justify inserting an arbitrary growing $m$ into a
fixed-parameter error term. The moment transition governed by
$n/(m+1)-\log m$ is a different simultaneous limit and should
remain separately indexed.

\paragraph{Height-one reductions and the literature.}
The Nielsen-polylogarithm work of Charlton, Gangl, and Radchenko
is a necessary part of the attribution for
Section~\ref{gorbit:sec:main}. In particular, the weight-seven
height-one reduction is an established functional equation.
Our presentation supplies a directly replayable shuffle
certificate, an explicit product correction, and a comparison
with the full binary quotient. The finite-direction formula
is stated in that full quotient, with its precise relation
space visible. Even when a formal obstruction is nonzero,
it must not be described as a proof that a complex number
cannot reduce.

\paragraph{Unresolved evaluations.}
The $S_6$ and revised $S_8$ candidates in the existing discovery
ledger remain conjectural. Neither the new Euler budget, a very
small residual, nor a formal obstruction to one reduction method
proves or disproves these special-value identities. This article
does not alter their evaluation status.

\subsection{The finite arithmetic proofs}

The exact programs use Python integers and
\src{fractions.Fraction}. No binary floating-point comparison
decides the sign of an interval in these proofs. The use of
Python is an implementation choice, not a trust claim about
all software: the written formulas specify finite statements
that can be replayed in another exact-arithmetic environment.
These are computer-assisted arithmetic proofs, not
proof-assistant formalizations.

\begin{longtable}{@{}>{\raggedright\arraybackslash}p{0.26\textwidth}>{\raggedright\arraybackslash}p{0.40\textwidth}>{\raggedright\arraybackslash}p{0.27\textwidth}@{}}
\toprule
Verifier & Exact task & Logical role\\
\midrule
\endhead
\src{verify_euler.py}
& Thirty-one rational axis enclosures, the rational global
budget, selected rational-order enclosures, and corrupted
candidate controls.
& Proves the finite inequalities used in
Proposition~\ref{prop:rational-budget}.\\[5pt]
\src{verify_depth_orbits.py}
& Exact Lyndon-bracket pairings, word-rank calculations,
the weight-seven product expansion, and the weight-eight
annihilator.
& Finite proof certificates for the explicit identities;
finite corroboration of the all-weight rank theorem.\\[5pt]
\src{certify_radial_degeneracy.py}
& Outward intervals for the simultaneous-zero rectangle,
the endpoint signs, the Jacobian, and the sixth coefficient.
& Supplies the finite sign claims needed for the analytic
existence and unfolding proof.\\[5pt]
\src{certify_two_extrema.py}
& Traps the actual implicit zero; bounds the infinite
series and its mixed derivatives; proves three signs
and uniform strict concavity.
& Proves the two-extremum statement for the exact
rational parameter pair.\\
\bottomrule
\end{longtable}
\addtocounter{table}{-1}

The Euler verifier also checks $6120$ rational kernel instances.
That grid is a useful regression check; the all-$N$ kernel theorem
is proved by the probability-mixture argument, not by extrapolation
from the grid. Its two deliberately false controls are a proposed
global constant $1$ and the changed finite value
$E_2(1,1)=-1/7$. The first is rejected by a certified analytic
enclosure; the second is rejected by exact rational evaluation,
which gives $-1/8$.

The word-rank checks through weight ten cover $25$ weight-depth
cases for each of the six-map and full-linear spans, and $34$
additional finite-direction cases. They are deliberately computed
from the expanded free-Lie words rather than by reusing the
dimension formula as the answer. The height-one membership checks
also test the asserted threshold on individual words. The
weight-eight obstruction is checked on $222$ transformed seed
words and $1792$ shuffle products. The weight-seven product
correction contains $864$ collected monomials; expanding their
shuffle products reproduces the residual word combination
exactly.

The first radial certificate uses a $2^{-192}$ outward grid.
Its logarithms are bounded by an atanh series with an explicit
geometric remainder, and exponentials by a positive Taylor sum
and a geometric tail. The second certificate range-reduces
logarithms before applying the same principle. The programs
record both the raw rational intervals and the simpler rational
enclosures printed in the theorems. Very small positive or
negative quantities are therefore proved to have their asserted
signs; their size is not inferred from decimal precision alone.

\subsection{Numerical diagnostics and independent review}

The optional numerical programs require \src{mpmath}, and the
asymptotic symbolic checks additionally use \src{sympy}.
They serve a different purpose from the exact programs.
The complete weight-seven functional identity, including
endpoint constants and products, was checked by direct disk
series at $z=1/4$ and $z=1/3$ with $80$ working decimal digits.
The discrepancies were approximately $3.7\cdot10^{-80}$ and
$2.1\cdot10^{-80}$. The proof comes from the exact word
identity and the branch argument; these numerical values
check their implementation.

The proposed sharp rate for $C_N-1$ is accompanied by $210$-digit
axis calculations and a few positive-$a$ checks. The latter replace
$g_{a,b}$ by $E_{512}$, with analytic contribution to the scaled
truncation error at most $(5/4)2^{N-512}$. That proved truncation
bound is separate from floating-point error. Neither these values
nor the numerical axis stationary points certify a global maximum.

For the reflected coefficients, the diagnostic data use $180$
working digits, compare two cases with independent Cauchy
quadrature, and repeat the six direct scaled-moment integrals
at $215$ digits. The repeated moment ratios differed by less
than $1.5\cdot10^{-108}$ in those cases. These are high-precision
agreements, not rigorous interval error bounds.
The accompanying tables retain examples outside the visible
asymptotic regime, including large normalized deviations at
moderate indices. Suppressing those rows would obscure the
meaning of a uniform theorem whose constants depend on a
fixed range parameter.

\begin{table}[tbp]
\centering
\small
\begin{tabular}{@{}rrrll@{}}
\toprule
$N$ & $n$ & $m$ & Observed normalized remainder
& First-correction prediction\\
\midrule
24 & 30 & 5  & $0.5127180890556$ & $0.5116978569285$\\
24 & 30 & 15 & $2.1726463253744$ & $0.5111559565783$\\
48 & 61 & 7  & $0.5094463157920$ & $0.5091282364092$\\
48 & 61 & 21 & $0.6787471547306$ & $0.5088627052376$\\
96 & 121 & 10 & $0.5027543017622$ & $0.5026340767530$\\
96 & 121 & 29 & $0.5028409944568$ & $0.5025085929532$\\
\bottomrule
\end{tabular}
\caption{Diagnostic values for the left side of
\eqref{moderate:eq:half}, compared with
$1/2+(3/2-\delta_\Gamma m^2/N-\eta)/(4N)$.
All six rows are retained, including the substantial
preasymptotic deviations at larger $m$ and smaller $N$.
The observed values are high-precision quadrature
results, not interval enclosures.}
\label{tab:half-diagnostics}
\end{table}

The four main arguments received separate mathematical
cross-checks during preparation. In particular, the endpoint
corrections in the Nielsen identity, the coefficient-dual
normalization in the binary-form calculation, the signed-cut
normalization in the half-term theorem, and the complete
implicit-derivative certificate for the explicit radial example
were checked independently of their initial derivations.
This record describes the scope of review; it is not a
substitute for external peer review.

\subsection{Reproduction and integration}

The archive includes a single entry point,
\src{code/reproduce.py}, which runs the exact verifiers without
requiring access to the original repository. Optional switches
replay the numerical diagnostics and regenerate the figures.
The figures are explanatory numerical plots; their captions
identify the exact statements that supply proof.
Run \src{latexmk -pdf article.tex} from the package root to
rebuild the document using the supplied section and figure files.

The provenance directory records the pinned commit, the
relative source paths used in the audit, and their SHA-256
digests. The integration guide maps each section and certificate
to the appropriate manuscript chapter. It also records the
small preamble requirements and label prefixes, so that the
standalone article can be integrated without depending on an
unstated local macro package. No changes were applied to the
source repository.
